\section{Introduction}
\label{sec:introduction}
Browsing the internet has become an integral part of everyday modern life. The internet is often subdivided into connected information networks that humans navigate using their internal models of the world. Understanding how humans navigate through such information networks gives insight into the connectedness of topics from the viewpoint of users. With this knowledge of semantic relations, the curators of these networks can organise their network nodes in a way that allows for more efficient and intuitive user navigation.

Wikipedia is a well-known example of a virtual information network that can be represented as a directed graph with articles as nodes and links between articles as the edges. It is partially structured by the semantic relatedness of the articles, and therefore acts as an approximation of humans' semantic world models. However, links between articles are added at the discretion of the article's author. While these links are often strongly semantically related, the amount of links and the degree of relatedness between linked articles can vary wildly. This makes it difficult to derive semantic relations from graph structure alone.

Wikispeedia is a game where the goal is for users to navigate from a source to a target article in a small version of Wikipedia (cite stuff here). Data collected from this game offers insight into the semantic relatedness of articles without direct connections, as well as the degree of degree of relatedness (the semantic distance) between articles. West et al. (cite wikispeedia article) demonstrated that semantic distance between concepts can be inferred from this game data.

In the interest of efficient navigation, it is worthwhile to investigate the relationship between the semantic distances of concepts and navigation speed. Understanding how semantic distance affects the time taken to move from a source to a target article offers oppurtunites to improve graph organisation. For example, rather than only having links added by authors, Wikipedia could provide additional links to concepts that minimises navigation time. This intuition can also be applied to the organisation of other semantic information networks. 

We aim to investigate this relationship by analysing Wikispeedia game data. However, due to the graph structure of Wikipedia, semantic distance is not the only factor influencing the duration of games. Factors such as the in-degree of an article (the amount of directed edges that have that article as the end node) also impact navigation difficulty. Defining an index of difficulty for Wikispeedia games that accounts for semantic distance as well as other confounding factors can be an initial step in investigating this relationship.

